% sections/introduction.tex

\section{Introduction}
\label{sec:intro}

Large language models (LLMs) are now widely deployed as automated judges---scoring RLHF outputs~\citep{ouyang2022training}, ranking chatbot responses~\citep{zheng2023judging}, and evaluating generation quality---but they make systematic errors, exhibit positional biases, and vary verdicts based on evaluation-protocol wording, all while expressing high confidence~\citep{li2024llmsasjudges,wang2023faireval,wei2024systematic}.

Prior work has established that LLM judges are overconfident and has begun to control this risk explicitly.
\citet{jung2025trust} showed that standard confidence measures (predictive probability, verbalized confidence) poorly predict whether a judge's verdict agrees with human annotators, and proposed \emph{Simulated Annotators} with cascaded selective evaluation to provide provable human-agreement guarantees; \citet{badshah2026scope} instead aggregate both response orderings into a permutation-invariant conformal score.
These methods each rely on a \emph{single} confidence source---simulated-annotator agreement or bidirectional preference probability---and do not exploit the broader observation that sensitivity across \emph{multiple} semantically equivalent evaluation protocols is itself an error signal.
We build directly on their selective-risk machinery (a held-out calibration threshold with an exact Clopper--Pearson bound) and do not claim it as a contribution; our contribution is the confidence signal that machinery is applied to. We are careful \emph{not} to over-claim its formal guarantees: because our acceptance threshold is chosen data-dependently on the calibration split, we treat the risk criterion as an empirical target validated by measured violation rates, not as a family-wise-corrected sequential test.

We investigate a complementary and deliberately asymmetric hypothesis: \textbf{a judgment that is unstable under semantically equivalent evaluation protocols is often unreliable, whereas stability alone is not sufficient evidence of correctness.}
Concretely, if an LLM judge produces different rankings when (i)~the rubric is paraphrased, (ii)~the response order is swapped, or (iii)~the judgment is re-sampled at nonzero temperature, the resulting verdict is less trustworthy.
The converse does not hold: a judge can be stably wrong, applying the same erroneous heuristic (e.g., a length or fluency preference) under every protocol variant~\citep{xu2026inside}.
We call this property \emph{protocol stability}, and we show that protocol \emph{instability} is a useful sufficient warning signal in competence regimes where the judge is at least moderately accurate---and we characterize precisely where it fails.

We organize our study around three questions: \emph{(RQ1)} does per-instance protocol stability predict judge error? \emph{(RQ2)} does fusing protocol-stability signals improve the risk--coverage tradeoff over standard confidence and single-signal baselines at a matched budget? \emph{(RQ3)} in which capability and task regimes do these signals help, and where do they break down?

Building on this insight, we introduce \textbf{CARE-Judge} (Calibrated, Abstaining, and Robust Evaluation), which (i)~collects multi-signal protocol-stability features---rubric perturbation, position-swap consistency, and stochastic self-consistency---for each pairwise judgment; (ii)~learns a correctness calibrator $\hat{p}(\text{correct}\mid\text{features})$ on a training split; (iii)~selects an acceptance threshold on a disjoint calibration split by an empirical prefix search certified with an exact Clopper--Pearson bound~\citep{jung2025trust,clopper1934use}; and (iv)~reports selective risk, coverage, and empirical risk-violation rates on a held-out test split.

We validate CARE-Judge across three judge models spanning a wide capability range (Qwen2.5-1.5B, DeepSeek-V4, GPT-5.5) and four benchmarks covering objective correctness (JudgeBench~\citep{tan2024judgebench}), reward-model preference (RewardBench~\citep{lambert2024rewardbench}), summarization preference (TL;DR~\citep{stiennon2020learning}), and open-ended human preference (LMArena~\citep{zheng2023judging}), totaling over 22{,}000 pairwise evaluations.
Our experiments answer the three research questions directly: protocol instability produces accuracy gaps of up to 27 points between stable and unstable subsets where the judge is at least moderately accurate but loses its predictive sign near chance (RQ1); at a matched six-call budget the fused calibrator attains the best or statistically-tied AUROC in 9 of 12 judge--dataset pairs, adding up to $+6.3$ AUROC points over the best single signal on the mid-capability DeepSeek-V4 (RQ2); and below a competence floor the signal \emph{inverts}, so that Qwen2.5-1.5B (49.4\% on JudgeBench) produces stable-but-wrong verdicts on which CARE-Judge correctly abstains (RQ3).

Our contributions are:
\begin{enumerate}
    \item \textbf{Competence-gated asymmetry (primary).} We establish---and quantify---that protocol instability predicts judge error only above a minimum judge-competence threshold, and that below it the signal \emph{inverts}: a near-random judge produces stable-but-wrong verdicts (Section~\ref{sec:exp:stablewrong}), a boundary we validate with a negative-control experiment showing that on a mid-capability judge non-equivalent rubrics flip $2.2\times$ more often than equivalent ones, whereas a near-random judge is insensitive to the distinction. To our knowledge this competence-gated asymmetry has not been characterized in prior per-instance judge-reliability work.
    \item \textbf{Decision-oriented multi-view fusion.} We fuse three families of protocol-perturbation instability---rubric paraphrase, response permutation, and stochastic resampling---into a \emph{single per-instance accept/abstain decision under an explicit risk budget}, and show it ranks correct versus incorrect verdicts better than standard confidence, single-signal baselines (including the SCOPE-style bidirectional signal), and ensemble aggregation at a matched inference budget (Section~\ref{sec:experiments}). Prompt- and order-sensitivity have been studied before as \emph{diagnostic} properties~\citep{hua2025flaw,bryan2025order,anghel2025diagnosing,huang2026prompt}, and individual consistency signals have been used to \emph{rank} or ensemble judges~\citep{li2025autoprompt,gupta2026diagnosing}; our contribution is neither observing protocol sensitivity nor a single consistency check, but combining all three perturbation families into one calibrated, risk-controlled acceptance rule (Section~\ref{sec:method}).
\end{enumerate}
We do not claim the risk-control machinery as novel; it is adopted from prior selective-judging work~\citep{jung2025trust,badshah2026scope}. Code and data will be released upon acceptance.
